\documentclass[11pt]{article}
\usepackage[a4paper,margin=25mm]{geometry}
\usepackage[T1]{fontenc}
\usepackage{amsmath}
\usepackage{booktabs}
\usepackage{graphicx}
\usepackage{longtable}
\usepackage{listings}
\usepackage{xcolor}
\usepackage{hyperref}

\hypersetup{
  colorlinks=true,
  linkcolor=blue,
  urlcolor=blue,
  citecolor=blue
}
\lstset{
  basicstyle=\ttfamily\small,
  breaklines=true,
  columns=fullflexible,
  keepspaces=true,
  showstringspaces=false
}
\title{macOS Computer Literacy\\Setup and Everyday Reference}
\author{UCC Physics Society}
\begin{document}
\maketitle

\tableofcontents
\section{Install VS Code or VSCodium}
Choose one editor. For VS Code, download the disk image and move
\texttt{Visual Studio Code.app} to Applications using the
\href{https://code.visualstudio.com/docs/setup/mac}{official macOS setup
guide}. For VSCodium, follow the
\href{https://vscodium.com/#install}{official installation options}; the
Homebrew command is available if Homebrew is already installed. Neither
Homebrew nor an administrator account is required for this workshop.

\section{Install the shared tools}
\begin{itemize}
\item Install Python from the
\href{https://www.python.org/downloads/macos/}{official Python macOS
downloads}. Use \texttt{python3}; do not modify Apple-managed system files.
\item Run \texttt{git --version}. If macOS offers to install Command Line
Tools, follow Apple's prompt; full Xcode is not required just to use Git.
\item For local LaTeX, install
\href{https://tug.org/mactex/}{MacTeX}. It is a large download; Overleaf is a
browser alternative.
\end{itemize}
These vendor pages provide current installation choices:
\href{https://git-scm.com/download/mac}{Git for macOS},
\href{https://brew.sh/}{Homebrew}, and
\href{https://tug.org/texlive/}{TeX Live}. Homebrew is optional, not a
prerequisite. After installing MacTeX, reopen Terminal before checking
\texttt{pdflatex --version} and \texttt{latexmk -v}.

\begin{lstlisting}
python3 --version
git --version
pwd
\end{lstlisting}


\newpage

\section{Terminal, projects, and paths}
Open Terminal from Spotlight. Recent macOS versions use \textbf{zsh} by
default. The terminal runs a shell to navigate folders and start programs.
Relative paths begin from the current working directory. Finder's
\texttt{Command+Shift+G} opens a path; \texttt{Command+Shift+.} toggles hidden
files. Treat filenames as case-sensitive for portable work.

\begin{lstlisting}
pendulum-lab/
  README.md
  requirements.txt
  data/
    raw_measurements.csv
  analysis/
    analysis.py
    analysis.ipynb
  figures/
  report/
    report.tex
\end{lstlisting}

\begin{longtable}{p{0.30\linewidth}p{0.63\linewidth}}
\toprule Task & Terminal command\\ \midrule
Current directory & \texttt{pwd}\\
List files, including hidden & \texttt{ls -la}\\
Move to the parent directory & \texttt{cd ..}\\
Go to your home folder & \texttt{cd \textasciitilde}\\
Enter a path with spaces & \texttt{cd "My Project"}\\
Copy / move a file & \texttt{cp source.csv copy.csv} / \texttt{mv copy.csv data.csv}\\
Find CSV files & \texttt{find . -name '*.csv'}\\
Search file text & \texttt{grep -n 'time' notes.txt}\\
Open the current folder in Finder & \texttt{open .}\\
Open the folder in the editor & \texttt{code .} or \texttt{codium .}\\
\bottomrule
\end{longtable}
Inspect paths before copying or moving files: an existing destination may be
overwritten. Prefer Finder's Trash when removing files. If the editor command
is unavailable, use File > Open Folder or install its shell command from the
editor's Command Palette.
See Apple's
\href{https://support.apple.com/guide/terminal/welcome/mac}{Terminal User
Guide} for opening Terminal, working with windows, and using commands.

\section{Python and project environments}
Create one virtual environment in the project directory containing
\texttt{requirements.txt}. Install project packages inside it:
\begin{lstlisting}
python3 -m venv .venv
source .venv/bin/activate
python -m pip install -r requirements.txt
python -m pip check
python -c "import sys; print(sys.executable)"
deactivate
\end{lstlisting}
The printed interpreter path should include \texttt{.venv}. A virtual
environment keeps project packages separate; it is not a security sandbox,
and it should not be committed or copied to another computer. See the
\href{https://docs.python.org/3/library/venv.html}{Python virtual-environment
documentation}. Activation applies only to the current terminal.

\section{Git and LaTeX}
Git works locally; GitHub hosts online copies and collaboration. Saving a file
is not the same as committing it, and a commit is not the same as pushing it.
For the shared commands and safe workflow, see the
\href{../Git/Git.pdf}{Git and GitHub guide}.

LaTeX source (\texttt{.tex}) is compiled into a PDF. Use Overleaf in a browser
or install MacTeX for local builds. From the repository root, build this
handout with:
\begin{lstlisting}
latexmk -cd -pdf -interaction=nonstopmode -halt-on-error MacOS/MaxOS.TEX
\end{lstlisting}
The shared \href{../LaTeX/LaTeX.pdf}{LaTeX guide} covers report structure,
equations, figures, and references; see also the
\href{https://www.overleaf.com/learn/latex/Learn_LaTeX_in_30_minutes}{Overleaf
LaTeX introduction} and
\href{https://www.latex-project.org/help/documentation/}{LaTeX Project
documentation}.

\section{Shortcuts and zsh configuration}
\begin{itemize}
\item \texttt{Ctrl+C} interrupts a terminal command; \texttt{Tab} completes
paths; \texttt{Up} recalls command history.
\item \texttt{Command+C/V} copies and pastes in standard macOS applications.
\item \texttt{Command+Space} opens Spotlight; \texttt{Command+Tab} switches
applications.
\end{itemize}

zsh reads \texttt{\textasciitilde/.zshrc} for interactive shells. To add
personal shortcuts, open the file in a text editor and append simple aliases:
\begin{lstlisting}
alias ll='ls -lah'
alias croot='cd ~'
\end{lstlisting}
Save the file, then load the changes in the current terminal:
\begin{lstlisting}
source ~/.zshrc
\end{lstlisting}
Back up and review your own file before editing it. Bash configuration is
different; do not copy a Linux \texttt{.bashrc} wholesale. See the
\href{https://zsh.sourceforge.io/Doc/Release/Files.html}{zsh startup-file
documentation}.

\section{Troubleshooting}
\begin{description}
\item[Python command missing:] Use \texttt{python3} and check that the
official installer completed.
\item[Editor command missing:] Use File > Open Folder or install the
editor's shell command, then restart Terminal.
\item[Python package installation is blocked:] Check that the project
environment is active.
\item[File not found:] Check \texttt{pwd}, filename case, and the relative path.
\item[TeX command missing:] Reopen Terminal after installing MacTeX; consult
its official installation instructions before changing paths.
\item[Application blocked:] Verify its source and university policy; do not
disable macOS security globally.
\end{description}

\noindent For Python, fitting, and data-analysis exercises, use the
\href{../Example_Code/Exersize.ipynb}{shared exercise notebook}.
\end{document}
